\documentclass[10pt]{extarticle}
\usepackage[a4paper,margin=1.5cm]{geometry}
\usepackage[T1]{fontenc}
\usepackage{amsmath,amssymb}
\usepackage{graphicx}
\usepackage[export]{adjustbox}
\usepackage{multicol}
\usepackage{enumitem}
\usepackage{booktabs,array,tabularx}
\setlength{\columnsep}{0.7cm}
\setlength{\parindent}{0pt}
\setlength{\parskip}{3pt}
\pagestyle{empty}
\makeatletter
\newcommand{\qfigdef}[4]{\@namedef{qf@file@#1}{#2}\@namedef{qf@w@#1}{#3}\@namedef{qf@wide@#1}{#4}}
\newcommand{\qf@check}[1]{\@ifundefined{qf@file@#1}{\PackageError{score2}{Unknown figure id #1}{}}{}}
\newcommand{\qf@img}[4][]{\includegraphics[width=#2\dimexpr\@nameuse{qf@w@#4}pt\relax,max width=#3,max height=0.45\textheight,#1]{\@nameuse{qf@file@#4}}}
\newif\ifqf@full
\newif\ifqf@boxed
\newsavebox{\qf@box}
\newcommand{\qf@open}{\global\qf@fulltrue\par\global\qf@boxedtrue
  \setbox\qf@box\vbox\bgroup\hsize\textwidth\linewidth\textwidth}
\newcommand{\qf@emit}[1]{\par\egroup\global\qf@boxedfalse
  \ifdim\dimexpr\ht\qf@box+\dp\qf@box\relax<#1\textheight\noindent\usebox{\qf@box}\par\else\unvbox\qf@box\fi}
\newcommand{\qfullsplit}{\ifqf@boxed\qf@emit{0.85}\fi}
\newcommand{\qf@close}{\ifqf@boxed\qf@emit{0.45}\fi\global\qf@fullfalse}
\newcommand{\qfullbegin}{\end{multicols}\qf@open}
\newcommand{\qfullend}{\qf@close\begin{multicols}{2}}
\newcommand{\qfullstart}{\qf@open}
\newcommand{\qfulldone}{\qf@close}
\newcommand{\qfig}[2][1]{\qf@check{#2}%
  \ifqf@full\qf@blk{#1}{#2}\else
  \ifnum\@nameuse{qf@wide@#2}=1 \expandafter\@firstoftwo\else\expandafter\@secondoftwo\fi
  {\end{multicols}\begin{center}\qf@img{#1}{\textwidth}{#2}\end{center}\begin{multicols}{2}}%
  {\qf@blk{#1}{#2}}\fi}
\newcommand{\qf@blk}[2]{\par\smallskip{\centering\qf@img{#1}{\linewidth}{#2}\par}\smallskip}
\newcommand{\qfigi}[2][1]{\qf@check{#2}\qf@img[valign=c]{#1}{\linewidth}{#2}}
\makeatother
\newcommand{\papersec}[1]{\par\bigskip{\centering\Large\bfseries PAPER-#1\par}\medskip}
\newcommand{\partsec}[2]{\par\medskip\noindent{\large\bfseries PART-#1: #2}\par\smallskip}
\newcommand{\sect}[1]{\par\smallskip\noindent\textbf{#1}\par}
\newcommand{\qn}[1]{\par\smallskip\noindent\textbf{#1.}~}
\newcommand{\paraq}[2]{\par\smallskip\noindent\textbf{Paragraph for Questions #1 and #2}\par}
\newenvironment{qoptions}{\begin{itemize}[leftmargin=1.8em,itemsep=1pt,topsep=2pt,label={}]}{\end{itemize}}
\newcommand{\opt}[2]{\item[(#1)] #2}
% figure registry (generated by score2_tools.py finalize)
\qfigdef{p1_phy_s1_q01_1}{JEE Advanced (2021 Series) - Test 02 (04 Feb 2021)-img/p1_phy_s1_q01_1.png}{211.8}{0}
\qfigdef{p1_phy_s1_q02_1}{JEE Advanced (2021 Series) - Test 02 (04 Feb 2021)-img/p1_phy_s1_q02_1.png}{153.2}{0}
\qfigdef{p1_phy_s1_q03_1}{JEE Advanced (2021 Series) - Test 02 (04 Feb 2021)-img/p1_phy_s1_q03_1.png}{70.0}{0}
\qfigdef{p1_phy_s1_q03_2}{JEE Advanced (2021 Series) - Test 02 (04 Feb 2021)-img/p1_phy_s1_q03_2.png}{141.4}{0}
\qfigdef{p1_phy_s1_q03_A}{JEE Advanced (2021 Series) - Test 02 (04 Feb 2021)-img/p1_phy_s1_q03_A.png}{106.3}{0}
\qfigdef{p1_phy_s1_q03_B}{JEE Advanced (2021 Series) - Test 02 (04 Feb 2021)-img/p1_phy_s1_q03_B.png}{128.5}{0}
\qfigdef{p1_phy_s1_q03_C}{JEE Advanced (2021 Series) - Test 02 (04 Feb 2021)-img/p1_phy_s1_q03_C.png}{116.7}{0}
\qfigdef{p1_phy_s1_q03_D}{JEE Advanced (2021 Series) - Test 02 (04 Feb 2021)-img/p1_phy_s1_q03_D.png}{138.7}{0}
\qfigdef{p1_phy_s1_q04_1}{JEE Advanced (2021 Series) - Test 02 (04 Feb 2021)-img/p1_phy_s1_q04_1.png}{123.5}{0}
\qfigdef{p1_phy_s1_q05_1}{JEE Advanced (2021 Series) - Test 02 (04 Feb 2021)-img/p1_phy_s1_q05_1.png}{168.2}{0}
\qfigdef{p1_phy_s1_q05_A}{JEE Advanced (2021 Series) - Test 02 (04 Feb 2021)-img/p1_phy_s1_q05_A.png}{12.2}{0}
\qfigdef{p1_phy_s1_q05_B}{JEE Advanced (2021 Series) - Test 02 (04 Feb 2021)-img/p1_phy_s1_q05_B.png}{13.5}{0}
\qfigdef{p1_phy_s1_q05_C}{JEE Advanced (2021 Series) - Test 02 (04 Feb 2021)-img/p1_phy_s1_q05_C.png}{12.9}{0}
\qfigdef{p1_phy_s1_q05_D}{JEE Advanced (2021 Series) - Test 02 (04 Feb 2021)-img/p1_phy_s1_q05_D.png}{12.9}{0}
\qfigdef{p1_phy_s1_q06_1}{JEE Advanced (2021 Series) - Test 02 (04 Feb 2021)-img/p1_phy_s1_q06_1.png}{228.0}{0}
\qfigdef{p1_phy_s2_q10_1}{JEE Advanced (2021 Series) - Test 02 (04 Feb 2021)-img/p1_phy_s2_q10_1.png}{148.4}{0}
\qfigdef{p1_phy_s2_q11_1}{JEE Advanced (2021 Series) - Test 02 (04 Feb 2021)-img/p1_phy_s2_q11_1.png}{219.3}{0}
\qfigdef{p1_phy_s2_q12_1}{JEE Advanced (2021 Series) - Test 02 (04 Feb 2021)-img/p1_phy_s2_q12_1.png}{117.4}{0}
\qfigdef{p1_phy_s3_q02_1}{JEE Advanced (2021 Series) - Test 02 (04 Feb 2021)-img/p1_phy_s3_q02_1.png}{96.0}{0}
\qfigdef{p1_phy_s3_q03_2}{JEE Advanced (2021 Series) - Test 02 (04 Feb 2021)-img/p1_phy_s3_q03_2.png}{232.4}{0}
\qfigdef{p1_chem_s1_q01_1}{JEE Advanced (2021 Series) - Test 02 (04 Feb 2021)-img/p1_chem_s1_q01_1.png}{124.8}{0}
\qfigdef{p1_chem_s1_q04_1}{JEE Advanced (2021 Series) - Test 02 (04 Feb 2021)-img/p1_chem_s1_q04_1.png}{395.5}{1}
\qfigdef{p1_chem_s1_q05_A}{JEE Advanced (2021 Series) - Test 02 (04 Feb 2021)-img/p1_chem_s1_q05_A.png}{78.7}{0}
\qfigdef{p1_chem_s1_q05_2}{JEE Advanced (2021 Series) - Test 02 (04 Feb 2021)-img/p1_chem_s1_q05_2.png}{59.8}{0}
\qfigdef{p1_chem_s1_q05_C}{JEE Advanced (2021 Series) - Test 02 (04 Feb 2021)-img/p1_chem_s1_q05_C.png}{78.7}{0}
\qfigdef{p1_chem_s1_q05_D}{JEE Advanced (2021 Series) - Test 02 (04 Feb 2021)-img/p1_chem_s1_q05_D.png}{79.2}{0}
\qfigdef{p1_chem_s1_q06_1}{JEE Advanced (2021 Series) - Test 02 (04 Feb 2021)-img/p1_chem_s1_q06_1.png}{55.5}{0}
\qfigdef{p1_chem_s1_q06_2}{JEE Advanced (2021 Series) - Test 02 (04 Feb 2021)-img/p1_chem_s1_q06_2.png}{109.3}{0}
\qfigdef{p1_chem_s1_q06_3}{JEE Advanced (2021 Series) - Test 02 (04 Feb 2021)-img/p1_chem_s1_q06_3.png}{106.0}{0}
\qfigdef{p1_chem_s1_q06_4}{JEE Advanced (2021 Series) - Test 02 (04 Feb 2021)-img/p1_chem_s1_q06_4.png}{109.3}{0}
\qfigdef{p1_chem_s1_q06_5}{JEE Advanced (2021 Series) - Test 02 (04 Feb 2021)-img/p1_chem_s1_q06_5.png}{106.0}{0}
\qfigdef{p1_chem_s2_q08_1}{JEE Advanced (2021 Series) - Test 02 (04 Feb 2021)-img/p1_chem_s2_q08_1.png}{364.2}{1}
\qfigdef{p1_chem_s2_q09_1}{JEE Advanced (2021 Series) - Test 02 (04 Feb 2021)-img/p1_chem_s2_q09_1.png}{384.0}{1}
\qfigdef{p1_chem_s2_q09_4}{JEE Advanced (2021 Series) - Test 02 (04 Feb 2021)-img/p1_chem_s2_q09_4.png}{125.3}{0}
\qfigdef{p1_chem_s2_q09_3}{JEE Advanced (2021 Series) - Test 02 (04 Feb 2021)-img/p1_chem_s2_q09_3.png}{131.8}{0}
\qfigdef{p1_chem_s2_q09_5}{JEE Advanced (2021 Series) - Test 02 (04 Feb 2021)-img/p1_chem_s2_q09_5.png}{128.7}{0}
\qfigdef{p1_chem_s2_q09_6}{JEE Advanced (2021 Series) - Test 02 (04 Feb 2021)-img/p1_chem_s2_q09_6.png}{125.5}{0}
\qfigdef{p1_chem_s2_q12_1}{JEE Advanced (2021 Series) - Test 02 (04 Feb 2021)-img/p1_chem_s2_q12_1.png}{91.6}{0}
\qfigdef{p1_chem_s3_q02_1}{JEE Advanced (2021 Series) - Test 02 (04 Feb 2021)-img/p1_chem_s3_q02_1.png}{139.5}{0}
\qfigdef{p1_chem_s3_q06_1}{JEE Advanced (2021 Series) - Test 02 (04 Feb 2021)-img/p1_chem_s3_q06_1.png}{307.6}{1}
\qfigdef{p2_phy_s1_q01_2}{JEE Advanced (2021 Series) - Test 02 (04 Feb 2021)-img/p2_phy_s1_q01_2.png}{89.5}{0}
\qfigdef{p2_phy_s1_q01_1}{JEE Advanced (2021 Series) - Test 02 (04 Feb 2021)-img/p2_phy_s1_q01_1.png}{93.8}{0}
\qfigdef{p2_phy_s1_q02_1}{JEE Advanced (2021 Series) - Test 02 (04 Feb 2021)-img/p2_phy_s1_q02_1.png}{245.2}{0}
\qfigdef{p2_phy_s1_q03_1}{JEE Advanced (2021 Series) - Test 02 (04 Feb 2021)-img/p2_phy_s1_q03_1.png}{51.3}{0}
\qfigdef{p2_phy_s2_q02_1}{JEE Advanced (2021 Series) - Test 02 (04 Feb 2021)-img/p2_phy_s2_q02_1.png}{245.2}{0}
\qfigdef{p2_phy_s2_q05_1}{JEE Advanced (2021 Series) - Test 02 (04 Feb 2021)-img/p2_phy_s2_q05_1.png}{153.2}{0}
\qfigdef{p2_phy_s2_q05_2}{JEE Advanced (2021 Series) - Test 02 (04 Feb 2021)-img/p2_phy_s2_q05_2.png}{138.7}{0}
\qfigdef{p2_phy_s3_q01_1}{JEE Advanced (2021 Series) - Test 02 (04 Feb 2021)-img/p2_phy_s3_q01_1.png}{221.9}{0}
\qfigdef{p2_phy_s3_q03_1}{JEE Advanced (2021 Series) - Test 02 (04 Feb 2021)-img/p2_phy_s3_q03_1.png}{158.4}{0}
\qfigdef{p2_phy_s3_q05_1}{JEE Advanced (2021 Series) - Test 02 (04 Feb 2021)-img/p2_phy_s3_q05_1.png}{135.9}{0}
\qfigdef{p2_phy_s3_q06_1}{JEE Advanced (2021 Series) - Test 02 (04 Feb 2021)-img/p2_phy_s3_q06_1.png}{76.2}{0}
\qfigdef{p2_chem_s1_q01_1}{JEE Advanced (2021 Series) - Test 02 (04 Feb 2021)-img/p2_chem_s1_q01_1.png}{247.4}{1}
\qfigdef{p2_chem_s1_q02_1}{JEE Advanced (2021 Series) - Test 02 (04 Feb 2021)-img/p2_chem_s1_q02_1.png}{189.2}{0}
\qfigdef{p2_chem_s1_q02_2}{JEE Advanced (2021 Series) - Test 02 (04 Feb 2021)-img/p2_chem_s1_q02_2.png}{53.0}{0}
\qfigdef{p2_chem_s1_q02_B}{JEE Advanced (2021 Series) - Test 02 (04 Feb 2021)-img/p2_chem_s1_q02_B.png}{52.2}{0}
\qfigdef{p2_chem_s1_q02_C}{JEE Advanced (2021 Series) - Test 02 (04 Feb 2021)-img/p2_chem_s1_q02_C.png}{59.4}{0}
\qfigdef{p2_chem_s1_q02_5}{JEE Advanced (2021 Series) - Test 02 (04 Feb 2021)-img/p2_chem_s1_q02_5.png}{68.5}{0}
\qfigdef{p2_chem_s1_q05_4}{JEE Advanced (2021 Series) - Test 02 (04 Feb 2021)-img/p2_chem_s1_q05_4.png}{31.0}{0}
\qfigdef{p2_chem_s1_q05_3}{JEE Advanced (2021 Series) - Test 02 (04 Feb 2021)-img/p2_chem_s1_q05_3.png}{69.4}{0}
\qfigdef{p2_chem_s1_q05_2}{JEE Advanced (2021 Series) - Test 02 (04 Feb 2021)-img/p2_chem_s1_q05_2.png}{83.5}{0}
\qfigdef{p2_chem_s1_q05_1}{JEE Advanced (2021 Series) - Test 02 (04 Feb 2021)-img/p2_chem_s1_q05_1.png}{83.4}{0}
\qfigdef{p2_chem_s1_q06_1}{JEE Advanced (2021 Series) - Test 02 (04 Feb 2021)-img/p2_chem_s1_q06_1.png}{351.9}{1}
\qfigdef{p2_chem_s1_q06_2}{JEE Advanced (2021 Series) - Test 02 (04 Feb 2021)-img/p2_chem_s1_q06_2.png}{33.4}{0}
\qfigdef{p2_chem_s1_q06_4}{JEE Advanced (2021 Series) - Test 02 (04 Feb 2021)-img/p2_chem_s1_q06_4.png}{54.4}{0}
\qfigdef{p2_chem_s1_q06_3}{JEE Advanced (2021 Series) - Test 02 (04 Feb 2021)-img/p2_chem_s1_q06_3.png}{68.5}{0}
\qfigdef{p2_chem_s2_q01_1}{JEE Advanced (2021 Series) - Test 02 (04 Feb 2021)-img/p2_chem_s2_q01_1.png}{111.7}{0}
\qfigdef{p2_chem_s2_q03_1}{JEE Advanced (2021 Series) - Test 02 (04 Feb 2021)-img/p2_chem_s2_q03_1.png}{99.5}{0}
\qfigdef{p2_chem_s2_q04_1}{JEE Advanced (2021 Series) - Test 02 (04 Feb 2021)-img/p2_chem_s2_q04_1.png}{153.2}{0}
\qfigdef{p2_chem_s3_q02_1}{JEE Advanced (2021 Series) - Test 02 (04 Feb 2021)-img/p2_chem_s3_q02_1.png}{87.7}{0}
\qfigdef{p2_chem_s3_q05_1}{JEE Advanced (2021 Series) - Test 02 (04 Feb 2021)-img/p2_chem_s3_q05_1.png}{159.9}{0}

\begin{document}
\begin{center}{\Large\bfseries JEE Advanced (2021 Series) --- Test 02 (04 Feb 2021)}\end{center}
\begin{multicols}{2}
\papersec{1}
\partsec{1}{PHYSICS}
\sect{Section-I(i)}
\qn{1} A football of radius $R$ is kept on a hole of radius $r$ ($r<R$) made on a plank kept horizontally. One end of the plank is now lifted so that it gets tilted making an angle $\theta$ from the horizontal as shown in the figure below. The maximum value of $\theta$ so that the football does not start rolling down the plank satisfies (figure is schematic and not drawn to scale) -
\qfig{p1_phy_s1_q01_1}
\begin{qoptions}
\opt{A}{$\sin\theta=\dfrac rR$}
\opt{B}{$\tan\theta=\dfrac rR$}
\opt{C}{$\sin\theta=\dfrac r{2R}$}
\opt{D}{$\cos\theta=\dfrac r{2R}$}
\end{qoptions}
\qn{2} A light disc made of aluminium (a nonmagnetic material) is kept horizontally and is free to rotate about its axis as shown in the figure. A strong magnet is held vertically at a point above the disc away from its axis. On revolving the magnet about the axis of the disc, the disc will (figure is schematic and not drawn to scale)-
\qfig{p1_phy_s1_q02_1}
\begin{qoptions}
\opt{A}{rotate in the direction opposite to the direction of magnet's motion}
\opt{B}{rotate in the same direction as the direction of magnet's motion}
\opt{C}{not rotate and its temperature will remain unchanged}
\opt{D}{not rotate but its temperature will slowly rise}
\end{qoptions}
\qn{3} A small roller of diameter 20 cm has an axle of diameter 10 cm (see figure below on the left). It is on a horizontal floor and a meter scale is positioned horizontally on its axle with one edge of the scale on top of the axle (see figure on the right). The scale is now pushed slowly on the axle so that it moves without slipping on the axle, and the roller starts rolling without slipping. After the roller has moved 50 cm, the position of the scale will look like (figures are schematic and not drawn to scale)-
\begin{center}
\qfigi{p1_phy_s1_q03_1}\hspace{1em}\qfigi{p1_phy_s1_q03_2}
\end{center}
\begin{qoptions}
\opt{A}{\qfigi{p1_phy_s1_q03_A}}
\opt{B}{\qfigi{p1_phy_s1_q03_B}}
\opt{C}{\qfigi{p1_phy_s1_q03_C}}
\opt{D}{\qfigi{p1_phy_s1_q03_D}}
\end{qoptions}
\qn{4} A circular coil of radius $R$ and $N$ turns has negligible resistance. As shown in the schematic figure, its two ends are connected to two wires and it is hanging by those wires with its plane being vertical. The wires are connected to a capacitor with charge $Q$ through a switch. The coil is in a horizontal uniform magnetic field $B_0$ parallel to the plane of the coil. When the switch is closed, the capacitor gets discharged through the coil in a very short time. By the time the capacitor is discharged fully, magnitude of the angular momentum gained by the coil will be (assume that the discharge time is so short that the coil has hardly rotated during this time)-
\qfig{p1_phy_s1_q04_1}
\begin{qoptions}
\opt{A}{$\dfrac\pi2NQB_0R^2$}
\opt{B}{$\pi NQB_0R^2$}
\opt{C}{$2\pi NQB_0R^2$}
\opt{D}{$4\pi NQB_0R^2$}
\end{qoptions}
\qn{5} A parallel beam of light strikes a piece of transparent glass having cross section as shown in the figure below. Correct shape of the emergent wavefront will be (figures are schematic and not drawn to scale)-
\qfig{p1_phy_s1_q05_1}
\begin{qoptions}
\opt{A}{\qfigi{p1_phy_s1_q05_A}}
\opt{B}{\qfigi{p1_phy_s1_q05_B}}
\opt{C}{\qfigi{p1_phy_s1_q05_C}}
\opt{D}{\qfigi{p1_phy_s1_q05_D}}
\end{qoptions}
\qn{6} An open-ended U-tube of uniform cross-sectional area contains water (density $10^3$ kg m$^{-3}$). Initially the water level stands at 0.29 m from the bottom in each arm. Kerosene oil (a water-immiscible liquid) of density 800 kg m$^{-3}$ is added to the left arm until its length is 0.1 m, as shown in the schematic figure below. The ratio $\left(\dfrac{h_1}{h_2}\right)$ of the heights of the liquid in the two arms is-
\qfig{p1_phy_s1_q06_1}
\begin{qoptions}
\opt{A}{$\dfrac{15}{14}$}
\opt{B}{$\dfrac{35}{33}$}
\opt{C}{$\dfrac76$}
\opt{D}{$\dfrac54$}
\end{qoptions}
\sect{Section-I(ii)}
\qn{7} A particle of mass $m$ moves in circular orbits with potential energy $V(r)=Fr$, where $F$ is a positive constant and $r$ is its distance from the origin. Its energies are calculated using the Bohr model. If the radius of the particle's orbit is denoted by $R$ and its speed and energy are denoted by $v$ and $E$, respectively, then for the $n^{\text{th}}$ orbit (here $h$ is the Planck's constant)-
\begin{qoptions}
\opt{A}{$R\propto n^{1/3}$ and $v\propto n^{2/3}$}
\opt{B}{$R\propto n^{2/3}$ and $v\propto n^{1/3}$}
\opt{C}{$E=\dfrac32\left(\dfrac{n^2h^2F^2}{4\pi^2m}\right)^{1/3}$}
\opt{D}{$E=2\left(\dfrac{n^2h^2F^2}{4\pi^2m}\right)^{1/3}$}
\end{qoptions}
\qn{8} The filament of a light bulb has surface area 64 mm$^2$. The filament can be considered as a black body at temperature 2500 K emitting radiation like a point source when viewed from far. At night the light bulb is observed from a distance of 100 m. Assume the pupil of the eyes of the observer to be circular with radius 3 mm. Then\\
(Take Stefan-Boltzmann constant $=5.67\times10^{-8}$ Wm$^{-2}$ K$^{-4}$, Wien's displacement constant $=2.90\times10^{-3}$ m-K, Planck's constant $=6.63\times10^{-34}$ Js, speed of light in vacuum $=3.00\times10^8$ ms$^{-1}$)-
\begin{qoptions}
\opt{A}{power radiated by the filament is in the range 642 W to 645 W}
\opt{B}{radiated power entering into one eye of the observer is in the range $3.15\times10^{-8}$ W to $3.25\times10^{-8}$ W}
\opt{C}{the wavelength corresponding to the maximum intensity of light is 1160 nm}
\opt{D}{taking the average wavelength of emitted radiation to be 1740 nm, the total number of photons entering per second into one eye of the observer is in the range $2.75\times10^{11}$ to $2.85\times10^{11}$}
\end{qoptions}
\qn{9} Sometimes it is convenient to construct a system of units so that all quantities can be expressed in terms of only one physical quantity. In one such system, dimensions of different quantities are given in terms of a quantity $X$ as follows: [position] $=[X^\alpha]$; [speed] $=[X^\beta]$; [acceleration] $=[X^p]$; [linear momentum] $=[X^q]$; [force] $=[X^r]$. Then -
\begin{qoptions}
\opt{A}{$\alpha+p=2\beta$}
\opt{B}{$p+q-r=\beta$}
\opt{C}{$p-q+r=\alpha$}
\opt{D}{$p+q+r=\beta$}
\end{qoptions}
\qn{10} A uniform electric field, $\vec E=-400\sqrt3\,\hat y$ NC$^{-1}$ is applied in a region. A charged particle of mass $m$ carrying positive charge $q$ is projected in this region with an initial speed of $2\sqrt{10}\times10^6$ ms$^{-1}$. This particle is aimed to hit a target T, which is 5 m away from its entry point into the field as shown schematically in the figure. Take $\dfrac qm=10^{10}$ Ckg$^{-1}$. Then:-
\qfig{p1_phy_s2_q10_1}
\begin{qoptions}
\opt{A}{the particle will hit T if projected at an angle $45^\circ$ from the horizontal}
\opt{B}{the particle will hit T if projected either at an angle $30^\circ$ or $60^\circ$ from the horizontal}
\opt{C}{time taken by the particle to hit T could be $\sqrt{\dfrac56}\ \mu$s as well as $\sqrt{\dfrac52}\ \mu$s}
\opt{D}{time taken by the particle to hit T is $\sqrt{\dfrac53}\ \mu$s}
\end{qoptions}
\qn{11} Shown in the figure is a semicircular metallic strip that has thickness $t$ and resistivity $\rho$. Its inner radius is $R_1$ and outer radius is $R_2$. If a voltage $V_0$ is applied between its two ends, a current $I$ flows in it. In addition, it is observed that a transverse voltage $\Delta V$ develops between its inner and outer surfaces due to purely kinetic effects of moving electrons (ignore any role of the magnetic field due to the current). Then (figure is schematic and not drawn to scale)-
\qfig{p1_phy_s2_q11_1}
\begin{qoptions}
\opt{A}{$I=\dfrac{V_0t}{\pi\rho}\ln\!\left(\dfrac{R_2}{R_1}\right)$}
\opt{B}{the outer surface is at a higher voltage than the inner surface}
\opt{C}{the outer surface is at a lower voltage than the inner surface}
\opt{D}{$\Delta V\propto I^2$}
\end{qoptions}
\qn{12} As shown schematically in the figure, two vessels contain water solutions (at temperature $T$) of potassium permanganate (KMnO$_4$) of different concentrations $n_1$ and $n_2$ ($n_1>n_2$) molecules per unit volume with $\Delta n=(n_1-n_2)\ll n_1$. When they are connected by a tube of small length $\ell$ and cross-sectional area $S$, KMnO$_4$ starts to diffuse from the left to the right vessel through the tube. Consider the collection of molecules to behave as dilute ideal gases and the difference in their partial pressure in the two vessels causing the diffusion. The speed $v$ of the molecules is limited by the viscous force $-\beta v$ on each molecule, where $\beta$ is a constant. Neglecting all terms of the order $(\Delta n)^2$, which of the following is/are correct? ($k_B$ is the Boltzmann constant)-
\qfig{p1_phy_s2_q12_1}
\begin{qoptions}
\opt{A}{the force causing the molecules to move across the tube is $\Delta nk_BTS$}
\opt{B}{force balance implies $n_1\beta v\ell=\Delta nk_BT$}
\opt{C}{total number of molecules going across the tube per sec is $\left(\dfrac{\Delta n}{\ell}\right)\left(\dfrac{k_BT}{\beta}\right)S$}
\opt{D}{rate of molecules getting transferred through the tube does not change with time}
\end{qoptions}
\sect{Section-II}
\qn{1} Put a uniform meter scale horizontally on your extended index fingers with the left one at 0.00 cm and the right one at 90.00 cm. When you attempt to move both the fingers slowly towards the center, initially only the left finger slips with respect to the scale and the right finger does not. After some distance, the left finger stops and the right one starts slipping. Then the right finger stops at a distance $x_R$ from the center (50.00 cm) of the scale and the left one starts slipping again. This happens because of the difference in the frictional forces on the two fingers. If the coefficients of static and dynamic friction between the fingers and the scale are 0.40 and 0.32, respectively, the value of $x_R$ (in cm) is \underline{\hspace{2cm}}.
\qn{2} When water is filled carefully in a glass, one can fill it to a height $h$ above the rim of the glass due to the surface tension of water. To calculate $h$ just before water starts flowing, model the shape of the water above the rim as a disc of thickness $h$ having semicircular edges, as shown schematically in the figure. When the pressure of water at the bottom of this disc exceeds what can be withstood due to the surface tension, the water surface breaks near the rim and water starts flowing from there. If the density of water, its surface tension and the acceleration due to gravity are $10^3$ kg m$^{-3}$, 0.07 Nm$^{-1}$ and 10 ms$^{-2}$, respectively, the value of $h$ (in mm) is \underline{\hspace{2cm}}.
\qfig{p1_phy_s3_q02_1}
\qn{3} One end of a spring of negligible unstretched length and spring constant $k$ is fixed at the origin (0,0). A point particle of mass $m$ carrying a positive charge $q$ is attached at its other end. The entire system is kept on a smooth horizontal surface. When a point dipole $\vec p$ pointing towards the charge $q$ is fixed at the origin, the spring gets stretched to a length $\ell$ and attains a new equilibrium position (see figure below). If the point mass is now displaced slightly by $\Delta\ell\ll\ell$ from its equilibrium position and released, it is found to oscillate at frequency $\dfrac1\delta\sqrt{\dfrac km}$. The value of $\delta$ is.
\qfig{p1_phy_s3_q03_2}
\qn{4} Consider one mole of helium gas enclosed in a container at initial pressure $P_1$ and volume $V_1$. It expands isothermally to volume $4V_1$. After this, the gas expands adiabatically and its volume becomes $32V_1$. The work done by the gas during isothermal and adiabatic expansion processes are $W_{\text{iso}}$ and $W_{\text{adia}}$, respectively. If the ratio $\dfrac{W_{\text{iso}}}{W_{\text{adia}}}=f\ln2$, then $f$ is \underline{\hspace{2cm}}.
\qn{5} A stationary tuning fork is in resonance with an air column in a pipe. If the tuning fork is moved with a speed of 2 ms$^{-1}$ in front of the open end of the pipe and parallel to it, the length of the pipe should be changed for the resonance to occur with the moving tuning fork. If the speed of sound in air is 320 ms$^{-1}$, the smallest value of the percentage change required in the length of the pipe is \underline{\hspace{2cm}}.
\qn{6} A circular disc of radius $R$ carries surface charge density $\sigma(r)=\sigma_0\left(1-\dfrac rR\right)$, where $\sigma_0$ is a constant and $r$ is the distance from the center of the disc. Electric flux through a large spherical surface that encloses the charged disc completely is $\phi_0$. Electric flux through another spherical surface of radius $\dfrac R4$ and concentric with the disc is $\phi$. Then the ratio $\dfrac{\phi_0}{\phi}$ is \underline{\hspace{2cm}}.
\partsec{2}{CHEMISTRY}
\sect{Section-I(i)}
\qn{1} If the distribution of molecular speeds of a gas is as per the figure shown below, then the ratio of the most probable, the average and the roots mean square speeds, respectively, is
\qfig{p1_chem_s1_q01_1}
\begin{qoptions}
\opt{A}{1 : 1 : 1}
\opt{B}{1 : 1 : 1.224}
\opt{C}{1 : 1.128 : 1.224}
\opt{D}{1 : 1.128 : 1}
\end{qoptions}
\qn{2} Which of the following liberates O$_2$ upon hydrolysis?
\begin{qoptions}
\opt{A}{Pb$_3$O$_4$}
\opt{B}{KO$_2$}
\opt{C}{Na$_2$O$_2$}
\opt{D}{Li$_2$O$_2$}
\end{qoptions}
\qn{3} A colorless aqueous solution contains nitrates of two metals, \textbf{X} and \textbf{Y}. When it was added to an aqueous solution of NaCl, a white precipitate was formed. This precipitate was found to be partly soluble in hot water to give a residue \textbf{P} and a solution \textbf{Q}. The residue P was soluble in aq. NH$_3$ and also in excess sodium thiosulfate. The hot solution \textbf{Q} gave a yellow precipitate with KI. The metals \textbf{X} and \textbf{Y}, respectively, are
\begin{qoptions}
\opt{A}{Ag and Pb}
\opt{B}{Ag and Cd}
\opt{C}{Cd and Pb}
\opt{D}{Cd and Zn}
\end{qoptions}
\qfullbegin
\qn{4} Newman projections \textbf{P}, \textbf{Q}, \textbf{R} and \textbf{S} are shown below :
\qfig{p1_chem_s1_q04_1}
\qfullsplit
Which one of the following options represents identical molecules ?
\begin{qoptions}
\opt{A}{P and Q}
\opt{B}{Q and S}
\opt{C}{Q and R}
\opt{D}{R and S}
\end{qoptions}
\qfullend
\qn{5} Which one of the following structures has the IUPAC name 3-ethynyl-2-hydroxy-4-methylhex-3-en-5-ynoic acid ?
\begin{qoptions}
\opt{A}{\qfigi{p1_chem_s1_q05_A}}
\opt{B}{\qfigi{p1_chem_s1_q05_2}}
\opt{C}{\qfigi{p1_chem_s1_q05_C}}
\opt{D}{\qfigi{p1_chem_s1_q05_D}}
\end{qoptions}
\qn{6} The Fischer projection of D-erythrose is shown below.
\begin{center}\qfigi{p1_chem_s1_q06_1}\\ D-Erythrose\end{center}
D-Erythrose and its isomers are listed as \textbf{P}, \textbf{Q}, \textbf{R} and \textbf{S} in \textbf{Column-I}. Choose the correct relationship of \textbf{P}, \textbf{Q}, \textbf{R} and \textbf{S} with D-erythrose from \textbf{Column-II}.\par\nobreak
\begin{center}
\begin{tabular}{@{}l@{\hspace{4pt}}c@{\hspace{4pt}}l@{}}
\textbf{Column-I} & & \textbf{Column-II}\\[2pt]
\textbf{(P)} & \qfigi{p1_chem_s1_q06_2} & 1.\ Diastereomer\\[4pt]
\textbf{(Q)} & \qfigi{p1_chem_s1_q06_3} & 2.\ Identical\\[4pt]
\textbf{(R)} & \qfigi{p1_chem_s1_q06_4} & 3.\ Enantiomer\\[4pt]
\textbf{(S)} & \qfigi{p1_chem_s1_q06_5} &
\end{tabular}
\end{center}
\begin{qoptions}
\opt{A}{P$\to$2, Q$\to$3, R$\to$2, S$\to$2}
\opt{B}{P$\to$3, Q$\to$1, R$\to$1, S$\to$2}
\opt{C}{P$\to$2, Q$\to$1, R$\to$1, S$\to$3}
\opt{D}{P$\to$2, Q$\to$3, R$\to$3, S$\to$1}
\end{qoptions}
\sect{Section-I(ii)}
\qn{7} In thermodynamics, the P-V work done is given by
\[ w=-\int dV\,P_{\text{ext}}\]
For a system undergoing a particular process, the work done is ,
\[ w=-\int dV\left(\frac{RT}{V-b}-\frac{a}{V^2}\right)\]
This equation is applicable to a
\begin{qoptions}
\opt{A}{System that satisfies the van der Waals equation of state.}
\opt{B}{Process that is reversible and isothermal.}
\opt{C}{Process that is reversible and adiabatic.}
\opt{D}{Process that is irreversible and at constant isothermal.}
\end{qoptions}
\qfullbegin
\qn{8} With respect to the compounds \textbf{I-V}, choose the correct statement(s).
\qfig{p1_chem_s2_q08_1}
\qfullsplit
\begin{qoptions}
\opt{A}{The acidity of compound \textbf{I} is due to delocalization in the conjugate base.}
\opt{B}{The conjugate base of compound \textbf{IV} is aromatic.}
\opt{C}{Compound \textbf{II} becomes more acidic, when it has a -NO$_2$ substituent.}
\opt{D}{The acidity of compounds follows the order \textbf{I $>$ IV $>$ V $>$ II $>$ III.}}
\end{qoptions}
\qn{9} In the reaction scheme shown below \textbf{Q}, \textbf{R} and \textbf{S} are the major products.
\qfig{p1_chem_s2_q09_1}
\qfullsplit
The correct structure of
\begin{qoptions}
\opt{A}{\textbf{S} is\\ \qfigi{p1_chem_s2_q09_4}}
\opt{B}{\textbf{Q} is\\ \qfigi{p1_chem_s2_q09_3}}
\opt{C}{\textbf{R} is\\ \qfigi{p1_chem_s2_q09_5}}
\opt{D}{\textbf{S} is\\ \qfigi{p1_chem_s2_q09_6}}
\end{qoptions}
\qfullend
\qn{10} Choose the correct statement(s) among the following :
\begin{qoptions}
\opt{A}{[FeCl$_4$]$^-$ has tetrahedral geometry.}
\opt{B}{[Co(en)(NH$_3$)$_2$Cl$_2$]$^+$ has 2 geometrical isomers.}
\opt{C}{[FeCl$_4$]$^-$ has higher spin-only magnetic moment than [Co(en)(NH$_3$)$_2$Cl$_2$]$^+$.}
\opt{D}{The cobalt ion in [Co(en)(NH$_3$)$_2$Cl$_2$]$^+$ has sp$^3$d$^2$ hybridization.}
\end{qoptions}
\qn{11} With respect to hypochlorite, chlorate and perchlorate ions, choose the correct statement(s).
\begin{qoptions}
\opt{A}{The hypochlorite ion is the strongest conjugate base.}
\opt{B}{The molecular shape of only chlorate ion is influenced by the lone pair of electrons of Cl.}
\opt{C}{The hypochlorite and chlorate ions disproportionate to give rise to identical set of ions.}
\opt{D}{The hypochlorite ion oxidizes the sulfite ion.}
\end{qoptions}
\qn{12} The cubic unit cell structure of a compound containing cation M and anion X is shown below. When compared to the anion, the cation has smaller ionic radius. Choose the correct statement(s).
\qfig{p1_chem_s2_q12_1}
\begin{qoptions}
\opt{A}{The empirical formula of the compound is MX.}
\opt{B}{The cation M and anion X have different coordination geometries.}
\opt{C}{The ratio of M-X bond length to the cubic unit cell edge length is 0.866.}
\opt{D}{The ratio of the ionic radii of cation M to anion X is 0.414.}
\end{qoptions}
\sect{Section-II}
\qn{1} 5.00 mL of 0.10 M oxalic acid solution taken in a conical flask is titrated against NaOH from a burette using phenolphthalein indicator. The volume of NaOH required for the appearance of permanent faint pink color is tabulated below for five experiments. What is the concentration, in molarity, of the NaOH solution ?\par\nobreak
\begin{center}
\begin{tabular}{|c|c|}
\hline Exp. No. & Vol. of NaOH (mL)\\ \hline
1 & 12.5\\ 2 & 10.5\\ 3 & 9.0\\ 4 & 9.0\\ 5 & 9.0\\ \hline
\end{tabular}
\end{center}
\qn{2} Consider the reaction $\mathbf A\rightleftharpoons\mathbf B$ at 1000 K. At time $t'$, the temperature of the system was increased to 2000 K and the system was allowed to reach equilibrium. Throughout this experiment the partial pressure of $\mathbf A$ was maintained at 1 bar. Given below is the plot of the partial pressure of $\mathbf B$ with time. What is the ratio of the standard Gibbs energy of the reaction at 1000 K to that at 2000 K?
\qfig{p1_chem_s3_q02_1}
\qn{3} Consider a 70\% efficient hydrogen-oxygen fuel cell working under standard conditions at 1 bar and 298 K. Its cell reaction is
\[ \text{H}_2(g)+\tfrac12\text{O}_2(g)\to\text{H}_2\text{O}(\ell)\]
The work derived from the cell on the consumption of $1.0\times10^{-3}$mol of H$_2$(g) is used to compress 1.00 mol of a monoatomic ideal gas in a thermally insulted container. What is the change in the temperature (in K) of the ideal gas ?\\
The standard reduction potentials for the two half-cells are given below.\\
O$_2$(g) + 4H$^+$ (aq.) + 4e$^-$ $\to$ 2H$_2$O ($\ell$) , $E^o=1.23$ V,\\
2H$^+$ (aq.) + 2e$^-$ $\to$ H$_2$(g), $E^o=0.00$V.\\
Use $F=96500$ C mol$^{-1}$, $R=8.314$ J mol$^{-1}$K$^{-1}$
\qn{4} Aluminium reacts with sulfuric acid to form aluminium sulfate and hydrogen. What is the volume of hydrogen gas in liters (L) produced at 300 K and 1.0 atm pressure, when 5.4 g of aluminium and 50.0 mL of 5.0 M sulfuric acid are combined for the reaction ?\\
(Use molar mass of aluminium as 27.0 g mol$^{-1}$ , $R=0.082$ atm L mol$^{-1}$ K$^{-1}$)
\qn{5} ${}^{238}_{92}$U is known to undergo radioactive decay to form by ${}^{206}_{82}$Pb emitting alpha and beta particles. A rock initially contained $68\times10^{-6}$ g of ${}^{238}_{92}$U. If the number of alpha particles that it would emit during its radioactive decay of ${}^{238}_{92}$U to ${}^{206}_{82}$Pb in three half-lives is $Z\times10^{18}$, then what is the value of $Z$ ?
\qfullbegin
\qn{6} In the following reaction, compound \textbf{Q} is obtained from compound \textbf{P} via an ionic intermediate
\qfig{p1_chem_s3_q06_1}
\qfullsplit
What is the degree of unsaturation of Q ?
\qfullend
\partsec{3}{MATHEMATICS}
\sect{Section-I(i)}
\qn{1} Suppose $a$, $b$ denote the distinct real roots of the quadratic polynomial $x^2+20x-2020$ and suppose $c$, $d$ denote the distinct complex roots of the quadratic polynomial $x^2-20x+2020$. Then the value of
\[ ac(a-c)+ad(a-d)+bc(b-c)+bd(b-d)\]
is
\begin{qoptions}
\opt{A}{0}
\opt{B}{8000}
\opt{C}{8080}
\opt{D}{16000}
\end{qoptions}
\qn{2} If the function $f:\mathbb R\to\mathbb R$ is defined by $f(x)=|x|(x-\sin x)$, then which of the following statements is \textbf{TRUE}?
\begin{qoptions}
\opt{A}{$f$ is one-one, but \textbf{NOT} onto}
\opt{B}{$f$ is onto, but \textbf{NOT} one-one}
\opt{C}{$f$ is \textbf{BOTH} one-one and onto}
\opt{D}{$f$ is \textbf{NEITHER} one-one \textbf{NOR} onto}
\end{qoptions}
\qn{3} Let the functions $f:\mathbb R\to\mathbb R$ and $g:\mathbb R\to\mathbb R$ be defined by
\[ f(x)=e^{x-1}-e^{-|x-1|}\ \text{ and }\ g(x)=\frac12\left(e^{x-1}+e^{1-x}\right).\]
Then the area of the region in the first quadrant bounded by the curves $y=f(x)$, $y=g(x)$ and $x=0$ is
\begin{qoptions}
\opt{A}{$\left(2-\sqrt3\right)+\dfrac12\left(e-e^{-1}\right)$}
\opt{B}{$\left(2+\sqrt3\right)+\dfrac12\left(e-e^{-1}\right)$}
\opt{C}{$\left(2-\sqrt3\right)+\dfrac12\left(e+e^{-1}\right)$}
\opt{D}{$\left(2+\sqrt3\right)+\dfrac12\left(e+e^{-1}\right)$}
\end{qoptions}
\qn{4} Let $a$, $b$ and $\lambda$ be positive real numbers. Suppose $P$ is an end point of the latus rectum of the parabola $y^2=4\lambda x$, and suppose the ellipse $\dfrac{x^2}{a^2}+\dfrac{y^2}{b^2}=1$ passes through the point $P$. If the tangents to the parabola and the ellipse at the point $P$ are perpendicular to each other, then the eccentricity of the ellipse is
\begin{qoptions}
\opt{A}{$\dfrac1{\sqrt2}$}
\opt{B}{$\dfrac12$}
\opt{C}{$\dfrac13$}
\opt{D}{$\dfrac25$}
\end{qoptions}
\qn{5} Let $C_1$ and $C_2$ be two biased coins such that the probabilities of getting head in a single toss are $\dfrac23$ and $\dfrac13$, respectively. Suppose $\alpha$ is the number of heads that appear when $C_1$ is tossed twice, independently, and suppose $\beta$ is the number of heads that appear when $C_2$ is tossed twice, independently. Then the probability that the roots of the quadratic polynomial $x^2-\alpha x+\beta$ are real and equal, is
\begin{qoptions}
\opt{A}{$\dfrac{40}{81}$}
\opt{B}{$\dfrac{20}{81}$}
\opt{C}{$\dfrac12$}
\opt{D}{$\dfrac14$}
\end{qoptions}
\qn{6} Consider all rectangles lying in the region
\[\left\{(x,y)\in\mathbb R\times\mathbb R:0\le x\le\frac\pi2\ \text{and}\ 0\le y\le2\sin(2x)\right\}\]
and having one side on the $x$-axis. The area of the rectangle which has the maximum perimeter among all such rectangles, is
\begin{qoptions}
\opt{A}{$\dfrac{3\pi}2$}
\opt{B}{$\pi$}
\opt{C}{$\dfrac{\pi}{2\sqrt3}$}
\opt{D}{$\dfrac{\pi\sqrt3}2$}
\end{qoptions}
\sect{Section-I(ii)}
\qn{7} Let the function $f:\mathbb R\to\mathbb R$ be defined by $f(x)=x^3-x^2+(x-1)\sin x$ and let $g:\mathbb R\to\mathbb R$ be an arbitrary function. Let $fg:\mathbb R\to\mathbb R$ be the product function defined by $(fg)(x)=f(x)g(x)$. Then which of the following statements is/are TRUE ?
\begin{qoptions}
\opt{A}{If $g$ is continuous at $x=1$, then $fg$ is differentiable at $x=1$}
\opt{B}{If $fg$ is differentiable at $x=1$, then $g$ is continuous at $x=1$}
\opt{C}{If $g$ is differentiable at $x=1$, then $fg$ is differentiable at $x=1$}
\opt{D}{If $fg$ is differentiable at $x=1$, then $g$ is differentiable at $x=1$}
\end{qoptions}
\qn{8} Let $M$ be a $3\times3$ invertible matrix with real entries and let $I$ denote the $3\times3$ identity matrix. If $M^{-1}=\text{adj (adj }M)$, then which of the following statement is/are ALWAYS TRUE ?
\begin{qoptions}
\opt{A}{$M=I$}
\opt{B}{$\det M=1$}
\opt{C}{$M^2=I$}
\opt{D}{$(\text{adj }M)^2=I$}
\end{qoptions}
\qn{9} Let $S$ be the set of all complex numbers $z$ satisfying $|z^2+z+1|=1$. Then which of the following statements is/are TRUE ?
\begin{qoptions}
\opt{A}{$\left|z+\dfrac12\right|\le\dfrac12$ for all $z\in S$}
\opt{B}{$|z|\le2$ for all $z\in S$}
\opt{C}{$\left|z+\dfrac12\right|\ge\dfrac12$ for all $z\in S$}
\opt{D}{The set $S$ has exactly four elements}
\end{qoptions}
\qn{10} Let $x$, $y$ and $z$ be positive real numbers. Suppose $x$, $y$ and $z$ are lengths of the sides of a triangle opposite to its angles $X$, $Y$ and $Z$, respectively. If
\[\tan\frac X2+\tan\frac Z2=\frac{2y}{x+y+z},\]
then which of the following statements is/are TRUE?
\begin{qoptions}
\opt{A}{$2Y=X+Z$}
\opt{B}{$Y=X+Z$}
\opt{C}{$\tan\dfrac X2=\dfrac x{y+z}$}
\opt{D}{$x^2+z^2-y^2=xz$}
\end{qoptions}
\qn{11} Let $L_1$ and $L_2$ be the following straight lines.
\begin{gather*}L_1:\ \frac{x-1}{1}=\frac y{-1}=\frac{z-1}{3}\quad\text{and}\\ L_2:\ \frac{x-1}{-3}=\frac y{-1}=\frac{z-1}{1}\end{gather*}
Suppose the straight line
\[L:\ \frac{x-\alpha}{\ell}=\frac{y-1}{m}=\frac{z-\gamma}{-2}\]
lies in the plane containing $L_1$ and $L_2$, and passes through the point of intersection of $L_1$ and $L_2$. If the line $L$ bisects the acute angle between the lines $L_1$ and $L_2$, then which of the following statements is/are TRUE?
\begin{qoptions}
\opt{A}{$\alpha-\gamma=3$}
\opt{B}{$\ell+m=2$}
\opt{C}{$\alpha-\gamma=1$}
\opt{D}{$\ell+m=0$}
\end{qoptions}
\qn{12} Which of the following inequalities is/are TRUE?
\begin{qoptions}
\opt{A}{$\displaystyle\int_0^1x\cos x\,dx\ge\frac38$}
\opt{B}{$\displaystyle\int_0^1x\sin x\,dx\ge\frac3{10}$}
\opt{C}{$\displaystyle\int_0^1x^2\cos x\,dx\ge\frac12$}
\opt{D}{$\displaystyle\int_0^1x^2\sin x\,dx\ge\frac29$}
\end{qoptions}
\sect{Section-II}
\qn{1} Let $m$ be the minimum possible value of $\log_3\left(3^{y_1}+3^{y_2}+3^{y_3}\right)$, where $y_1,y_2,y_3$ are real numbers for which $y_1+y_2+y_3=9$. Let $M$ be the maximum possible value of $(\log_3x_1+\log_3x_2+\log_3x_3)$, where $x_1,x_2,x_3$ are positive real numbers for which $x_1+x_2+x_3=9$. Then the value of $\log_2(m^3)+\log_3(M^2)$ is \underline{\hspace{1.5cm}}.
\qn{2} Let $a_1,a_2,a_3,\ldots$ be a sequence of positive integers in arithmetic progression with common difference 2. Also, let $b_1,b_2,b_3,\ldots$ be a sequence of positive integers in geometric progression with common ratio 2. If $a_1=b_1=c$, then the number of all possible values of $c$, for which the equality
\[2(a_1+a_2+\ldots+a_n)=b_1+b_2+\ldots+b_n\]
holds for some positive integer $n$, is \underline{\hspace{1.5cm}}
\qn{3} Let $f:[0,2]\to\mathbb R$ be the function defined by
\[f(x)=(3-\sin(2\pi x))\sin\!\left(\pi x-\frac\pi4\right)-\sin\!\left(3\pi x+\frac\pi4\right)\]
If $\alpha,\beta\in[0,2]$ are such that $\{x\in[0,2]:f(x)\ge0\}=[\alpha,\beta]$, then the value of $\beta-\alpha$ is \underline{\hspace{1.5cm}}
\qn{4} In a triangle PQR, let $\vec a=\overrightarrow{QR}$, $\vec b=\overrightarrow{RP}$ and $\vec c=\overrightarrow{PQ}$. If $|\vec a|=3$, $|\vec b|=4$ and $\dfrac{\vec a\cdot(\vec c-\vec b)}{\vec c\cdot(\vec a-\vec b)}=\dfrac{|\vec a|}{|\vec a|+|\vec b|}$, then the value of $|\vec a\times\vec b|^2$ is \underline{\hspace{1.5cm}}
\qn{5} For a polynomial $g(x)$ with real coefficient, let $m_g$ denote the number of distinct real roots of $g(x)$. Suppose $S$ is the set of polynomials with real coefficients defined by\\
$S=\{(x^2-1)^2(a_0+a_1x+a_2x^2+a_3x^3):a_0,a_1,a_2,a_3\in\mathbb R\}$.\\
For a polynomial $f$, let $f'$ and $f''$ denote its first and second order derivatives, respectively. Then the minimum possible value of $(m_{f'}+m_{f''})$, where $f\in S$, is \underline{\hspace{1.5cm}}
\qn{6} Let $e$ denote the base of the natural logarithm. The value of the real number $a$ for which the right hand limit
\[\lim_{x\to0^+}\frac{(1-x)^{\frac1x}-e^{-1}}{x^a}\]
is equal to a nonzero real number, is \underline{\hspace{1.5cm}}.
\papersec{2}
\partsec{1}{PHYSICS}
\sect{Section-I}
\qn{1} A beaker of radius $r$ is filled with water (refractive index $\dfrac43$) up to a height $H$ as shown in the figure on the left. The beaker is kept on a horizontal table rotating with angular speed $\omega$. This makes the water surface curved so that the difference in the height of water level at the center and at the circumference of the beaker is $h$ ($h\ll H$, $h\ll r$), as shown in the figure on the right. Take this surface to be approximately spherical with a radius of curvature $R$. Which of the following is/are correct ? ($g$ is the acceleration due to gravity)
\begin{center}
\qfigi{p2_phy_s1_q01_2}\hspace{1em}\qfigi{p2_phy_s1_q01_1}
\end{center}
\begin{qoptions}
\opt{A}{$R=\dfrac{h^2+r^2}{2h}$}
\opt{B}{$R=\dfrac{3r^2}{2h}$}
\opt{C}{Apparent depth of the bottom of the beaker is close to $\dfrac{3H}{2}\left(1+\dfrac{\omega^2H}{2g}\right)^{-1}$}
\opt{D}{Apparent depth of the bottom of the beaker is close to $\dfrac{3H}{4}\left(1+\dfrac{\omega^2H}{4g}\right)^{-1}$}
\end{qoptions}
\qn{2} A student skates up a ramp that makes an angle $30^\circ$ with the horizontal. He/she starts (as shown in the figure) at the bottom of the ramp with speed $v_0$ and wants to turn around over a semicircular path $xyz$ of radius $R$ during which he/she reaches a maximum height $h$ (at point $y$) from the ground as shown in the figure. Assume that the energy loss is negligible and the force required for this turn at the highest point is provided by his/her weight only. Then ($g$ is the acceleration due to gravity)
\qfig{p2_phy_s1_q02_1}
\begin{qoptions}
\opt{A}{$v_0^2-2gh=\dfrac12 gR$}
\opt{B}{$v_0^2-2gh=\dfrac{\sqrt3}{2}gR$}
\opt{C}{the centripetal force required at points $x$ and $z$ is zero}
\opt{D}{the centripetal force required is maximum at points $x$ and $z$}
\end{qoptions}
\qn{3} A rod of mass $m$ and length $L$, pivoted at one of its ends, is hanging vertically. A bullet of the same mass moving at speed $v$ strikes the rod horizontally at a distance $x$ from its pivoted end and gets embedded in it. The combined system now rotates with angular speed $\omega$ about the pivot. The maximum angular speed $\omega_M$ is achieved for $x=x_M$. Then
\qfig{p2_phy_s1_q03_1}
\begin{qoptions}
\opt{A}{$\omega=\dfrac{3vx}{L^2+3x^2}$}
\opt{B}{$\omega=\dfrac{12vx}{L^2+12x^2}$}
\opt{C}{$x_M=\dfrac{L}{\sqrt3}$}
\opt{D}{$\omega_M=\dfrac{v}{2L}\sqrt3$}
\end{qoptions}
\qn{4} In an X-ray tube, electrons emitted from a filament (cathode) carrying current $I$ hit a target (anode) at a distance $d$ from the cathode. The target is kept at a potential $V$ higher than the cathode resulting in emission of continuous and characteristic X-rays. If the filament current $I$ is decreased to $\dfrac I2$, the potential difference $V$ is increased to $2V$, and the separation distance $d$ is reduced to $\dfrac d2$, then
\begin{qoptions}
\opt{A}{the cut-off wavelength will reduce to half, and the wavelengths of the characteristic X-rays will remain the same}
\opt{B}{the cut-off wavelength as well as the wavelengths of the characteristic X-rays will remain the same}
\opt{C}{the cut-off wavelength will reduce to half, and the intensities of all the X-rays will decrease}
\opt{D}{the cut-off wavelength will become two times larger, and the intensity of all the X-rays will decrease}
\end{qoptions}
\qn{5} Two identical non-conducting solid spheres of same mass and charge are suspended in air from a common point by two non-conducting, massless strings of same length. At equilibrium, the angle between the strings is $\alpha$. The spheres are now immersed in a dielectric liquid of density $800$ kg m$^{-3}$ and dielectric constant $21$. If the angle between the strings remains the same after the immersion, then
\begin{qoptions}
\opt{A}{electric force between the spheres remains unchanged}
\opt{B}{electric force between the spheres reduces}
\opt{C}{mass density of the spheres is $840$ kg m$^{-3}$}
\opt{D}{the tension in the strings holding the spheres remains unchanged}
\end{qoptions}
\qn{6} Starting at time $t=0$ from the origin with speed $1$ ms$^{-1}$, a particle follows a two-dimensional trajectory in the $x$-$y$ plane so that its coordinates are related by the equation $y=\dfrac{x^2}{2}$. The $x$ and $y$ components of its acceleration are denoted by $a_x$ and $a_y$, respectively. Then
\begin{qoptions}
\opt{A}{$a_x=1$ ms$^{-2}$ implies that when the particle is at the origin, $a_y=1$ ms$^{-2}$}
\opt{B}{$a_x=0$ implies $a_y=1$ ms$^{-2}$ at all times}
\opt{C}{at $t=0$, the particle's velocity points in the $x$-direction}
\opt{D}{$a_x=0$ implies that at $t=1$ s, the angle between the particle's velocity and the $x$ axis is $45^\circ$}
\end{qoptions}
\sect{Section-II}
\qn{1} A spherical bubble inside water has radius $R$. Take the pressure inside the bubble and the water pressure to be $p_0$. The bubble now gets compressed radially in an adiabatic manner so that its radius becomes $(R-a)$. For $a\ll R$ the magnitude of the work done in the process is given by $(4\pi p_0Ra^2)X$, where $X$ is a constant and $\gamma=\dfrac{C_P}{C_V}=\dfrac{41}{30}$. The value of $X$ is \underline{\hspace{2cm}}.
\qn{2} In the balanced condition, the values of the resistances of the four arms of a Wheatstone bridge are shown in the figure below. The resistance $R_3$ has temperature coefficient $0.0004\,^\circ\mathrm{C}^{-1}$. If the temperature of $R_3$ is increased by $100^\circ$C, the voltage developed between $S$ and $T$ will be \underline{\hspace{2cm}} volt.
\qfig{p2_phy_s2_q02_1}
\qn{3} Two capacitors with capacitance values $C_1=2000\pm10$ pF and $C_2=3000\pm15$ pF are connected in series. The voltage applied across this combination is $V=5.00\pm0.02$ V. The percentage error in the calculation of the energy stored in this combination of capacitors is \underline{\hspace{2cm}}.
\qn{4} A cubical solid aluminium (bulk modulus $=-V\dfrac{dP}{dV}=70$ GPa) block has an edge length of $1$ m on the surface of the earth. It is kept on the floor of a $5$ km deep ocean. Taking the average density of water and the acceleration due to gravity to be $10^3$ kg m$^{-3}$ and $10$ ms$^{-2}$, respectively, the change in the edge length of the block in mm is \underline{\hspace{2cm}}.
\qn{5} The inductors of two $LR$ circuits are placed next to each other, as shown in the figure. The values of the self-inductance of the inductors, resistances, mutual-inductance and applied voltages are specified in the given circuit. After both the switches are closed simultaneously, the total work done by the batteries against the induced EMF in the inductors by the time the currents reach their steady state values is \underline{\hspace{2cm}} mJ.
\begin{center}
\qfigi{p2_phy_s2_q05_1}\hspace{1em}\qfigi{p2_phy_s2_q05_2}
\end{center}
\qn{6} A container with $1$ kg of water in it is kept in sunlight, which causes the water to get warmer than the surroundings. The average energy per unit time per unit area received due to the sunlight is $700$ Wm$^{-2}$ and it is absorbed by the water over an effective area of $0.05$ m$^2$. Assuming that the heat loss from the water to the surroundings is governed by Newton's law of cooling, the difference (in $^\circ$C) in the temperature of water and the surroundings after a long time will be \underline{\hspace{2cm}}. (Ignore effect of the container, and take constant for Newton's law of cooling $=0.001$ s$^{-1}$, Heat capacity of water $=4200$ J kg$^{-1}$ K$^{-1}$)
\sect{Section-III}
\qn{1} A large square container with thin transparent vertical walls and filled with water (refractive index $\dfrac43$) is kept on a horizontal table. A student holds a thin straight wire vertically inside the water $12$ cm from one of its corners, as shown schematically in the figure. Looking at the wire from this corner, another student sees two images of the wire, located symmetrically on each side of the line of sight as shown. The separation (in cm) between these images is \underline{\hspace{2cm}}.
\qfig{p2_phy_s3_q01_1}
\qn{2} A train with cross-sectional area $S_t$ is moving with speed $v_t$ inside a long tunnel of cross-sectional area $S_0$ ($S_0=4S_t$). Assume that almost all the air (density $\rho$) in front of the train flows back between its sides and the walls of the tunnel. Also, the air flow with respect to the train is steady and laminar. Take the ambient pressure and that inside the train to be $p_0$. If the pressure in the region between the sides of the train and the tunnel walls is $p$, then $p_0-p=\dfrac{7}{2N}\rho v_t^2$. The value of $N$ is \underline{\hspace{2cm}}.
\qn{3} Two large circular discs separated by a distance of $0.01$ m are connected to a battery via a switch as shown in the figure. Charged oil drops of density $900$ kg m$^{-3}$ are released through a tiny hole at the center of the top disc. Once some oil drops achieve terminal velocity, the switch is closed to apply a voltage of $200$ V across the discs. As a result, an oil drop of radius $8\times10^{-7}$ m stops moving vertically and floats between the discs. The number of electrons present in this oil drop is \underline{\hspace{2cm}}. (neglect the buoyancy force, take acceleration due to gravity $=10$ ms$^{-2}$ and charge on an electron $(e)=1.6\times10^{-19}$ C)
\qfig{p2_phy_s3_q03_1}
\qn{4} A hot air balloon is carrying some passengers, and a few sandbags of mass $1$ kg each so that its total mass is $480$ kg. Its effective volume giving the balloon its buoyancy is $V$. The balloon is floating at an equilibrium height of $100$ m. When $N$ number of sandbags are thrown out, the balloon rises to a new equilibrium height close to $150$ m with its volume $V$ remaining unchanged. If the variation of the density of air with height $h$ from the ground is $\rho(h)=\rho_0e^{-h/h_0}$, where $\rho_0=1.25$ kg m$^{-3}$ and $h_0=6000$ m, the value of $N$ is \underline{\hspace{2cm}}.
\qn{5} A point charge $q$ of mass $m$ is suspended vertically by a string of length $\ell$. A point dipole of dipole moment $\vec p$ is now brought towards $q$ from infinity so that the charge moves away. The final equilibrium position of the system including the direction of the dipole, the angles and distances is shown in the figure below. If the work done in bringing the dipole to this position is $N\times(mgh)$, where $g$ is the acceleration due to gravity, then the value of $N$ is \underline{\hspace{2cm}}. (Note that for three coplanar forces keeping a point mass in equilibrium, $\dfrac{F}{\sin\theta}$ is the same for all forces, where $F$ is any one of the forces and $\theta$ is the angle between the other two forces)
\qfig{p2_phy_s3_q05_1}
\qn{6} A thermally isolated cylindrical closed vessel of height $8$ m is kept vertically. It is divided into two equal parts by a diathermic (perfect thermal conductor) frictionless partition of mass $8.3$ kg. Thus the partition is held initially at a distance of $4$ m from the top, as shown in the schematic figure below. Each of the two parts of the vessel contains $0.1$ mole of an ideal gas at temperature $300$ K. The partition is now released and moves without any gas leaking from one part of the vessel to the other. When equilibrium is reached, the distance of the partition from the top (in m) will be \underline{\hspace{2cm}}\\
(take the acceleration due to gravity $=10$ ms$^{-2}$ and the universal gas constant $=8.3$ J mol$^{-1}$K$^{-1}$).
\qfig{p2_phy_s3_q06_1}
\qfullbegin
\partsec{2}{CHEMISTRY}
\sect{Section-I}
\qn{1} In an experiment, $m$ grams of a compound \textbf{X} (gas/liquid/solid) taken in a container is loaded in a balance as shown in figure \textbf{I} below. In the presence of a magnetic field, the pan with \textbf{X} is either deflected upwards (figure \textbf{II}), or deflected downwards (figure \textbf{III}), depending on the compound \textbf{X}. Identify the correct statement(s)
\qfig{p2_chem_s1_q01_1}
\qfullsplit
\begin{qoptions}
\opt{A}{If \textbf{X} is H$_2$O($\ell$), deflection of the pan is upwards.}
\opt{B}{If \textbf{X} is K$_4$[Fe(CN)$_6$](s), deflection of the pan is upwards.}
\opt{C}{If \textbf{X} is O$_2$(g), deflection of the pan is downwards.}
\opt{D}{If \textbf{X} is C$_6$H$_6$($\ell$), deflection of the pan is downwards.}
\end{qoptions}
\qfullend
\qn{2} Which of the following plots is(are) correct for the given reaction? ([P]$_0$ is the initial concentration of \textbf{P})
\qfig{p2_chem_s1_q02_1}
\begin{qoptions}
\opt{A}{\qfigi{p2_chem_s1_q02_2}}
\opt{B}{\qfigi{p2_chem_s1_q02_B}}
\opt{C}{\qfigi{p2_chem_s1_q02_C}}
\opt{D}{\qfigi{p2_chem_s1_q02_5}}
\end{qoptions}
\qn{3} Which among the following statement(s) is(are) true for the extraction of aluminium from bauxite?
\begin{qoptions}
\opt{A}{Hydrated Al$_2$O$_3$ precipitates, when CO$_2$ is bubbled through a solution of sodium aluminate.}
\opt{B}{Addition of Na$_3$AlF$_6$ lowers the melting point of alumina.}
\opt{C}{CO$_2$ is evolved at the anode during electrolysis.}
\opt{D}{The cathode is a steel vessel with a lining of carbon.}
\end{qoptions}
\qn{4} Choose the correct statement(s) among the following.
\begin{qoptions}
\opt{A}{SnCl$_2\cdot$2H$_2$O is a reducing agent.}
\opt{B}{SnO$_2$ reacts with KOH to form K$_2$[Sn(OH)$_6$].}
\opt{C}{A solution of PbCl$_2$ in HCl contains Pb$^{2+}$ and Cl$^-$ ions.}
\opt{D}{The reaction of Pb$_3$O$_4$ with hot dilute nitric acid to give PbO$_2$ is a redox reaction.}
\end{qoptions}
\qn{5} Consider the following four compounds \textbf{I}, \textbf{II}, \textbf{III}, and \textbf{IV}.
\begin{center}
\qfigi{p2_chem_s1_q05_4}\hspace{0.5em}\qfigi{p2_chem_s1_q05_3}\hspace{0.5em}\qfigi{p2_chem_s1_q05_2}\hspace{0.5em}\qfigi{p2_chem_s1_q05_1}
\end{center}
Choose the correct statement(s).
\begin{qoptions}
\opt{A}{The order of basicity is \textbf{II} $>$ \textbf{I} $>$ \textbf{III} $>$ \textbf{IV}.}
\opt{B}{The magnitude of p$K_\mathrm{b}$ difference between \textbf{I} and \textbf{II} is more than that between \textbf{III} and \textbf{IV}.}
\opt{C}{Resonance effect is more in \textbf{III} than in \textbf{IV}.}
\opt{D}{Steric effect makes compound \textbf{IV} more basic than \textbf{III}.}
\end{qoptions}
\qfullbegin
\qn{6} Consider the following transformations of a compound \textbf{P}.
\qfig{p2_chem_s1_q06_1}
\qfullsplit
Choose the correct option(s).
\begin{qoptions}
\opt{A}{\textbf{P} is \qfigi{p2_chem_s1_q06_2}}
\opt{B}{\textbf{X} is Pd-C/quinoline/H$_2$}
\opt{C}{\textbf{P} is \qfigi{p2_chem_s1_q06_4}}
\opt{D}{\textbf{R} is \qfigi{p2_chem_s1_q06_3}}
\end{qoptions}
\qfullend
\sect{Section-II}
\qn{1} A solution of $0.1$ M weak base (B) is titrated with $0.1$ M of a strong acid (HA). The variation of pH of the solution with the volume of HA added is shown in the figure below. What is the p$K_\mathrm{b}$ of the base? The neutralization reaction is given by B + HA $\rightarrow$ BH$^+$ + A$^-$.
\qfig{p2_chem_s2_q01_1}
\qn{2} Liquids \textbf{A} and \textbf{B} form ideal solution for all compositions of \textbf{A} and \textbf{B} at $25^\circ$C. Two such solutions with $0.25$ and $0.50$ mole fractions of \textbf{A} have the total vapor pressures of $0.3$ and $0.4$ bar, respectively. What is the vapor pressure of pure liquid \textbf{B} in bar?
\qn{3} The figure below is the plot of potential energy versus internuclear distance ($d$) of H$_2$ molecule in the electronic ground state. What is the value of the net potential energy $E_0$ (as indicated in the figure) in kJ mol$^{-1}$, for $d=d_0$ at which the electron-electron repulsion and the nucleus-nucleus repulsion energies are absent? As reference, the potential energy of H atom is taken as zero when its electron and the nucleus are infinitely far apart.\\
Use Avogadro constant as $6.023\times10^{23}$ mol$^{-1}$.
\qfig{p2_chem_s2_q03_1}
\qn{4} Consider the reaction sequence from \textbf{P} to \textbf{Q} shown below. The overall yield of the major product \textbf{Q} from \textbf{P} is 75\%. What is the amount in grams of \textbf{Q} obtained from 9.3 mL of \textbf{P}? (Use density of \textbf{P} $=1.00$ g mL$^{-1}$, Molar mass of C $=12.0$, H $=1.0$, O $=16.0$ and N $=14.0$ g mol$^{-1}$)
\qfig{p2_chem_s2_q04_1}
\qn{5} Tin is obtained from cassiterite by reduction with coke. Use the data given below to determine the minimum temperature (in K) at which the reduction of cassiterite by coke would take place.\\
At 298 K: $\Delta_f H^0(\mathrm{SnO_2(s)})=-581.0$ kJ mol$^{-1}$, $\Delta_f H^0(\mathrm{CO_2(g)})=-394.0$ kJ mol$^{-1}$\\
$S^0(\mathrm{SnO_2(s)})=56.0$ J K$^{-1}$ mol$^{-1}$, $S^0(\mathrm{Sn(s)})=52.0$ J K$^{-1}$ mol$^{-1}$,\\
$S^0(\mathrm{C(s)})=6.0$ J K$^{-1}$ mol$^{-1}$, $S^0(\mathrm{CO_2(g)})=210.0$ J K$^{-1}$ mol$^{-1}$.\\
Assume that the enthalpies and the entropies are temperature independent.
\qn{6} An acidified solution of $0.05$ M Zn$^{2+}$ is saturated with $0.1$ M H$_2$S. What is the minimum molar concentration (M) of H$^+$ required to prevent the precipitation of ZnS ?\\
Use $K_\mathrm{sp}$(ZnS) $=1.25\times10^{-22}$ and overall dissociation constant of H$_2$S, $K_\mathrm{NET}=K_1K_2=1\times10^{-21}$
\sect{Section-III}
\qn{1} The $1^\mathrm{st}$, $2^\mathrm{nd}$ and the $3^\mathrm{rd}$ ionization enthalpies, $I_1$, $I_2$ and $I_3$, of four atoms with atomic numbers $n$, $n+1$, $n+2$ and $n+3$, where $n<10$, are tabulated below. What is the value of $n$?
\begin{center}
\begin{tabular}{|c|c|c|c|}
\hline
Atomic & \multicolumn{3}{c|}{Ionization Enthalpy (kJ/mol)}\\
\cline{2-4}
number & $I_1$ & $I_2$ & $I_3$\\
\hline
$n$ & 1681 & 3374 & 6050\\
$n+1$ & 2081 & 3952 & 6122\\
$n+2$ & 496 & 4562 & 6910\\
$n+3$ & 738 & 1451 & 7733\\
\hline
\end{tabular}
\end{center}
\qn{2} Consider the following compounds in the liquid form : O$_2$, HF, H$_2$O, NH$_3$, H$_2$O$_2$, CCl$_4$, CHCl$_3$, C$_6$H$_6$, C$_6$H$_5$Cl.\\
When a charged comb is brought near their flowing stream, how many of them show deflection as per the following figure?
\qfig{p2_chem_s3_q02_1}
\qn{3} In the chemical reaction between stoichiometric quantities of KMnO$_4$ and KI in weakly basic solution, what is the number of moles of I$_2$ released for 4 moles of KMnO$_4$ consumed ?
\qn{4} An acidified solution of potassium chromate was layered with an equal volume of amyl alcohol. When it was shaken after the addition of 1 mL of 3\% H$_2$O$_2$, a blue alcohol layer was obtained. The blue color is due to the formation of a chromium (VI) compound `\textbf{X}'. What is the number of oxygen atoms bonded to chromium through only single bonds in a molecule of \textbf{X}?
\qn{5} The structure of a peptide is given below.
\qfig{p2_chem_s3_q05_1}
If the absolute values of the net charge of the peptide at pH $=2$, pH $=6$, and pH $=11$ are $|z_1|$, $|z_2|$ and $|z_3|$, respectively, then what is $|z_1|+|z_2|+|z_3|$ ?
\qn{6} An organic compound (C$_8$H$_{10}$O$_2$) rotates plane-polarized light. It produces pink color with neutral FeCl$_3$ solution. What is the total number of all the possible isomers for this compound?
\columnbreak
\partsec{3}{MATHEMATICS}
\sect{Section-I}
\qn{1} Let $b$ be a nonzero real number. Suppose $f:\mathbb R\to\mathbb R$ is a differentiable function such that $f(0)=1$. If the derivative $f'$ of $f$ satisfies the equation $f'(x)=\dfrac{f(x)}{b^2+x^2}$ for all $x\in\mathbb R$, then which of the following statements is/are TRUE ?
\begin{qoptions}
\opt{A}{If $b>0$, then $f$ is an increasing function}
\opt{B}{If $b<0$, then $f$ is a decreasing function}
\opt{C}{$f(x)\,f(-x)=1$ for all $x\in\mathbb R$}
\opt{D}{$f(x)-f(-x)=0$ for all $x\in\mathbb R$}
\end{qoptions}
\qn{2} Let $a$ and $b$ be positive real numbers such that $a>1$ and $b<a$. Let $P$ be a point in the first quadrant that lies on the hyperbola $\dfrac{x^2}{a^2}-\dfrac{y^2}{b^2}=1$. Suppose the tangent to the hyperbola at $P$ passes through the point $(1,0)$, and suppose the normal to the hyperbola at $P$ cuts off equal intercepts on the coordinate axes. Let $\Delta$ denote the area of the triangle formed by the tangent at $P$, the normal at $P$ and the $x$-axis. If $e$ denotes the eccentricity of the hyperbola, then which of the following statements is/are TRUE ?
\begin{qoptions}
\opt{A}{$1<e<\sqrt2$}
\opt{B}{$\sqrt2<e<2$}
\opt{C}{$\Delta=a^4$}
\opt{D}{$\Delta=b^4$}
\end{qoptions}
\qn{3} Let $f:\mathbb R\to\mathbb R$ and $g:\mathbb R\to\mathbb R$ be functions satisfying $f(x+y)=f(x)+f(y)+f(x)f(y)$ and $f(x)=xg(x)$ for all $x,y\in\mathbb R$. If $\displaystyle\lim_{x\to0}g(x)=1$, then which of the following statements is/are TRUE?
\begin{qoptions}
\opt{A}{$f$ is differentiable at every $x\in\mathbb R$}
\opt{B}{If $g(0)=1$, then $g$ is differentiable at every $x\in\mathbb R$}
\opt{C}{The derivative $f'(1)$ is equal to $1$}
\opt{D}{The derivative $f'(0)$ is equal to $1$}
\end{qoptions}
\qn{4} Let $\alpha,\beta,\gamma,\delta$ be real numbers such that $\alpha^2+\beta^2+\gamma^2\ne0$ and $\alpha+\gamma=1$. Suppose the point $(3,2,-1)$ is the mirror image of the point $(1,0,-1)$ with respect to the plane $\alpha x+\beta y+\gamma z=\delta$. Then which of the following statements is/are TRUE ?
\begin{qoptions}
\opt{A}{$\alpha+\beta=2$}
\opt{B}{$\delta-\gamma=3$}
\opt{C}{$\delta+\beta=4$}
\opt{D}{$\alpha+\beta+\gamma=\delta$}
\end{qoptions}
\qn{5} Let $a$ and $b$ be positive real numbers. Suppose $\overrightarrow{PQ}=a\hat i+b\hat j$ and $\overrightarrow{PS}=a\hat i-b\hat j$ are adjacent sides of a parallelogram $PQRS$. Let $\vec u$ and $\vec v$ be the projection vectors of $\vec w=\hat i+\hat j$ along $\overrightarrow{PQ}$ and $\overrightarrow{PS}$, respectively. If $|\vec u|+|\vec v|=|\vec w|$ and if the area of the parallelogram $PQRS$ is $8$, then which of the following statements is/are TRUE ?
\begin{qoptions}
\opt{A}{$a+b=4$}
\opt{B}{$a-b=2$}
\opt{C}{The length of the diagonal $PR$ of the parallelogram $PQRS$ is $4$}
\opt{D}{$\vec w$ is an angle bisector of the vectors $\overrightarrow{PQ}$ and $\overrightarrow{PS}$}
\end{qoptions}
\qn{6} For non-negative integers $s$ and $r$, let
\[\binom sr=\begin{cases}\dfrac{s!}{r!(s-r)!}&\text{if }r\le s,\\[1ex] 0&\text{if }r>s.\end{cases}\]
For positive integers $m$ and $n$, let
\[g(m,n)=\sum_{p=0}^{m+n}\frac{f(m,n,p)}{\binom{n+p}{p}}\]
where for any nonnegative integer $p$,
\[f(m,n,p)=\sum_{i=0}^{p}\binom mi\binom{n+i}{p}\binom{p+n}{p-i}.\]
Then which of the following statements is/are TRUE?
\begin{qoptions}
\opt{A}{$g(m,n)=g(n,m)$ for all positive integers $m,n$}
\opt{B}{$g(m,n+1)=g(m+1,n)$ for all positive integers $m,n$}
\opt{C}{$g(2m,2n)=2g(m,n)$ for all positive integers $m,n$}
\opt{D}{$g(2m,2n)=(g(m,n))^2$ for all positive integers $m,n$}
\end{qoptions}
\sect{Section-II}
\qn{1} An engineer is required to visit a factory for exactly four days during the first 15 days of every month and it is mandatory that no two visits take place on consecutive days. Then the number of all possible ways in which such visits to the factory can be made by the engineer during 1-15 June 2021 is \underline{\hspace{2cm}}.
\qn{2} In a hotel, four rooms are available. Six persons are to be accommodated in these four rooms in such a way that each of these rooms contains at least one person and at most two persons. Then the number of all possible ways in which this can be done is \underline{\hspace{2cm}}.
\qn{3} Two fair dice, each with faces numbered $1,2,3,4,5$ and $6$, are rolled together and the sum of the numbers on the faces is observed. This process is repeated till the sum is either a prime number or a perfect square. Suppose the sum turns out to be a perfect square before it turns out to be a prime number. If $p$ is the probability that this perfect square is an odd number, then the value of $14p$ is \underline{\hspace{2cm}}.
\qn{4} Let the function $f:[0,1]\to\mathbb R$ be defined by $f(x)=\dfrac{4^x}{4^x+2}$. Then the value of $f\!\left(\dfrac1{40}\right)+f\!\left(\dfrac2{40}\right)+f\!\left(\dfrac3{40}\right)+\cdots+f\!\left(\dfrac{39}{40}\right)-f\!\left(\dfrac12\right)$ is \underline{\hspace{2cm}}.
\qn{5} Let $f:\mathbb R\to\mathbb R$ be a differentiable function such that its derivative $f'$ is continuous and $f(\pi)=-6$.\\
If $F:[0,\pi]\to\mathbb R$ is defined by $F(x)=\displaystyle\int_0^x f(t)\,dt$, and if
\[\int_0^\pi\bigl(f'(x)+F(x)\bigr)\cos x\,dx=2,\]
then the value of $f(0)$ is \underline{\hspace{2cm}}.
\qn{6} Let the function $f:(0,\pi)\to\mathbb R$ be defined by
\[f(\theta)=(\sin\theta+\cos\theta)^2+(\sin\theta-\cos\theta)^4\]
Suppose the function $f$ has a local minimum at $\theta$ precisely when $\theta\in\{\lambda_1\pi,\dots,\lambda_r\pi\}$, where $0<\lambda_1<\cdots<\lambda_r<1$. Then the value of $\lambda_1+\cdots+\lambda_r$ is \underline{\hspace{2cm}}.
\sect{Section-III}
\qn{1} For a complex number $z$, let $\mathrm{Re}(z)$ denote the real part of $z$. Let $S$ be the set of all complex numbers $z$ satisfying $z^4-|z|^4=4iz^2$, where $i=\sqrt{-1}$. Then the minimum possible value of $|z_1-z_2|^2$, where $z_1,z_2\in S$ with $\mathrm{Re}(z_1)>0$ and $\mathrm{Re}(z_2)<0$, is \underline{\hspace{2cm}}.
\qn{2} The probability that a missile hits a target successfully is $0.75$. In order to destroy the target completely, at least three successful hits are required. Then the minimum number of missiles that have to be fired so that the probability of completely destroying the target is NOT less than $0.95$, is \underline{\hspace{2cm}}.
\qn{3} Let $O$ be the centre of the circle $x^2+y^2=r^2$, where $r>\dfrac{\sqrt5}{2}$. Suppose $PQ$ is a chord of this circle and the equation of the line passing through $P$ and $Q$ is $2x+4y=5$. If the centre of the circumcircle of the triangle $OPQ$ lies on the line $x+2y=4$, then the value of $r$ is \underline{\hspace{2cm}}.
\qn{4} The trace of a square matrix is defined to be the sum of its diagonal entries. If $A$ is a $2\times2$ matrix such that the trace of $A$ is $3$ and the trace of $A^3$ is $-18$, then the value of the determinant of $A$ is \underline{\hspace{2cm}}.
\qn{5} Let the functions $f:(-1,1)\to\mathbb R$ and $g:(-1,1)\to(-1,1)$ be defined by $f(x)=|2x-1|+|2x+1|$ and $g(x)=x-[x]$, where $[x]$ denotes the greatest integer less than or equal to $x$. Let $f\circ g:(-1,1)\to\mathbb R$ be the composite function defined by $(f\circ g)(x)=f(g(x))$. Suppose $c$ is the number of points in the interval $(-1,1)$ at which $f\circ g$ is NOT continuous, and suppose $d$ is the number of points in the interval $(-1,1)$ at which $f\circ g$ is NOT differentiable. Then the value of $c+d$ is \underline{\hspace{2cm}}.
\qn{6} The value of the limit
\[\lim_{x\to\frac\pi2}\frac{4\sqrt2(\sin3x+\sin x)}{\begin{gathered}\left(2\sin2x\sin\tfrac{3x}2+\cos\tfrac{5x}2\right)\\-\left(\sqrt2+\sqrt2\cos2x+\cos\tfrac{3x}2\right)\end{gathered}}\]
is \underline{\hspace{2cm}}.
\end{multicols}
\end{document}
